% TOP_PROBLEM: {"id": 3254, "title": "Cycles in Graphs of Large Chromatic Number", "category_id": 3, "category": {"id": 3, "name": "graph_theory", "display_name": "Graph Theory", "description": "Problems involving graphs, networks, and their properties.", "slug": "graph-theory", "order_index": 3, "created_at": "2026-07-31T15:26:25.671Z"}, "status": "open", "problem_number": "OPG-60001", "set_id": 12, "statusQualification": "The missing hypothesis k>=3 is restored from the published formulation and explicitly explained."}
% FIRST_POSTED: 2026-10-02
% ATTEMPT_STATUS: unresolved
% UnsolvedMath ID 3254; source catalogue code OPG-60001
% !TeX program = lualatex
% Research attempt by GPT-6 Astra Ultra; no claim of a new complete solution.
\documentclass[11pt,a4paper]{article}
\usepackage[margin=25mm]{geometry}
\usepackage{fontspec}
\setmainfont{lmroman10-regular.otf}[BoldFont=lmroman10-bold.otf,
 ItalicFont=lmroman10-italic.otf,BoldItalicFont=lmroman10-bolditalic.otf]
\usepackage{amsmath,amssymb,mathtools}
\usepackage{unicode-math}
\setmathfont{latinmodern-math.otf}
\AtBeginDocument{\let\setminus\smallsetminus}
\usepackage{xurl,hyperref}
\hypersetup{colorlinks=true,urlcolor=blue}
\setcounter{secnumdepth}{0}
\setlength{\parindent}{0pt}
\setlength{\parskip}{5pt}
\setlength{\emergencystretch}{3em}
\begin{document}
\section*{Cycles in Graphs of Large Chromatic Number}\label{3254}
{\small UnsolvedMath ID 3254 · OPG-60001 · Graph Theory · OpenGarden}
\subsection{Definitions and mathematical statement}
Let $G$ be a finite simple undirected graph and let $k\ge3$ be an integer.
The chromatic number $\chi(G)$ is the least number of colours in a proper
vertex colouring. A cycle means a simple cycle, with no repeated vertex
except the equality of its initial and final vertex. Its length is its
number of edges. Count cycles as distinct subgraphs, without distinguishing
a starting vertex or a direction of traversal; chords between vertices
of a cycle are permitted.

Let $c_{0,k}(G)$ count the cycles whose lengths are positive multiples of
$k$. The intended conjecture is
\[
 \chi(G)>k\quad\Longrightarrow\quad
 c_{0,k}(G)\ge\frac{(k+1)(k-1)!}{2}.
\]
The bound counts cycles individually, with no requirement of edge or
vertex disjointness. In $K_{k+1}$ the only possible positive multiple of
$k$ among cycle lengths is $k$. Choosing the omitted vertex and a cyclic
ordering of the others gives exactly $(k+1)(k-1)!/2$ such cycles, explaining
the proposed extremal quantity.

The restriction $k\ge3$ is essential and is stated explicitly in the
primary formulation [S3]. It is missing from the imported one-sentence
statement. Without it the assertion fails: $K_3$ is more than
2-chromatic but has no even cycle, and a nontrivial tree is more than
1-chromatic but has no cycle. The present expansion restores this
published parameter range rather than treating those trivial failures
as the intended problem.
\subsection{Short English statement}
Does every graph needing more than k colours contain at least as many cycles of length divisible by k as the complete graph on k+1 vertices, for k at least three?
\subsection{Research attempt}
\paragraph{Scope and process.}
This is a GPT-6 Astra Ultra research attempt dated 2 October 2026, using
primary-source browsing and elementary mathematical checks. The canonical
definition above is retained. Results below are restricted results or
reductions, not a claim of a new solution to the entire catalogue question.
No formal proof assistant or external mathematical referee checked this text.


\paragraph{Literature and counting convention.}
The source attributes this conjecture to Brewster, McGuinness, Moore and
Noel. The primary paper [S3], checked again for this attempt, proves the
full case $k=3$. It also proves additional results under criticality and
minimum-degree hypotheses. We do not assume that a critical graph has
minimum degree exactly $k$: only a lower bound follows automatically.
The following argument establishes the required count for an infinite
family having chromatic number $k+1$ but no $K_{k+1}$. Each cycle is
counted up to rotation and reversal, as in the definition.

\paragraph{A critical reduction and the clique benchmark.}
If $\chi(G)>k$, choose an induced subgraph $H$ minimal by vertex deletion
with this property. Deleting any vertex leaves a $k$-colourable graph,
and restoring it with a fresh colour shows $\chi(H)=k+1$.
Every vertex of $H$ has degree at least $k$: otherwise a $k$-colouring
of its deletion extends by an unused colour. Since every cycle of $H$
is also a cycle of $G$, it would suffice to prove the target for these
critical graphs. This reduction does not bound their orders.

If $G$ contains $K_{k+1}$, each choice of $k$ vertices supports
$(k-1)!/2$ distinct cycles of length $k$. Different choices have different
vertex sets, giving exactly the proposed lower bound from that clique.
The reduction therefore leaves particular interest in critical graphs
without a clique of order $k+1$.

\paragraph{A nonclique family meeting the required count.}
Fix an odd integer $\ell\ge\max(5,k+1)$ and form the join
\[
                 H=K_{k-2}+C_\ell.
\]
A join retains the two graphs and adds every edge between them. Colours
used on its two sides must be disjoint. Since an odd cycle needs three
colours, $\chi(H)=(k-2)+3=k+1$. Its clique number is $k$, because a
cycle of length at least five has clique number two; in particular this
family is not covered by the preceding $K_{k+1}$ argument.

For each edge $uv$ of the cycle, take $u,v$ together with all $k-2$
vertices on the complete side. These $k$ vertices induce a $K_k$, which
has $(k-1)!/2$ Hamilton cycles. The justification is to fix one starting
vertex, order the remaining $k-1$ vertices, and identify the two
orientations of each cycle. The identification has no fixed point
because $k\ge3$. Cycles obtained from different rim edges have different
vertex sets, so none is counted twice. Consequently
\[
 c_{0,k}(H)\ \ge\ \ell\frac{(k-1)!}{2}
       \ \ge\ \frac{(k+1)(k-1)!}{2}.
\]
There may be many further cycles; the argument needs only these
length-$k$ witnesses. For $k=3$ this specializes to the odd wheel, whose
rim edges give distinct triangles through the centre. For every $k$ it
works for arbitrarily large odd $\ell$.

\paragraph{Verification and limits of a subdivision approach.}
The finite certificate [R1] constructs the joins with
$(k,\ell)=(3,5),(4,5),(5,7),(6,7)$ and enumerates their simple
length-$k$ cycles, identifying rotations and reversal. It checks the
lower count and the absence of a clique of order $k+1$. This verifies
the counting conventions on small representatives; the displayed proof
establishes the entire family.

Trying to replace the actual clique by a subdivided clique loses the
argument: subdivision changes cycle lengths and their residues modulo
$k$. It also need not preserve the relevant chromatic obstruction. For
example, subdividing every edge once makes any graph bipartite, by
separating original vertices from subdivision vertices. Thus the
existence of a topological clique alone gives neither the required
chromatic reduction nor the correct divisibility count. Likewise, the
critical reduction only yields $\delta(H)\ge k$, not equality; a theorem
requiring a vertex of degree $k$ needs an additional argument.

\paragraph{Final status and first remaining gap.}
\textbf{UNRESOLVED.} The known full $k=3$ theorem is acknowledged, and
an explicit clique-free-at-order-$k+1$ join family is proved to satisfy
the bound for every $k\ge3$. For general $k\ge4$, no argument here
produces the required number of distinct divisible-length cycles in
arbitrary critical graphs. Controlling that count outside the join
construction is the first remaining gap.

\subsection{Sources}
\begin{itemize}
\item[S1] ULAM.AI and the UnsolvedMath Contributors, supplied problems.json, record 3254 (OPG-60001), OpenGarden collection. Catalogue identity and source transcription; CC BY 4.0.
\url{https://huggingface.co/datasets/ulamai/UnsolvedMath}.
\item[S2] Open Problem Garden, Cycles in Graphs of Large Chromatic Number. Original problem and discussion; the mathematical formulation below is expanded from the supplied source record.
\url{http://www.openproblemgarden.org/op/cycles_in_graphs_of_large_chromatic_number}.
\item[S3] S. Kim and M. E. Picollelli, 4-Chromatic Graphs Have At Least 4 Cycles of Length 0 mod 3 (2023), abstract explicitly states k at least 3.
\url{https://arxiv.org/abs/2312.05945}.

\item[R1] Reproducible finite checks, using Python's standard library:
\path{scripts/check_attempt_batch03_colouring_2026_10_02.py}, section 3254.
The source [S3] was checked on 2 October 2026 for the stated literature result.

\end{itemize}
\end{document}
